\section{存活几何与硬件转化}
\label{sec:cost-hypothesis}

执行器观测每个问题首次通过验证器的时刻。这些迭代编号定义存活曲线，并由此确定可移除工作。后端宽度成本曲线进一步决定工作量减少中有多少能够转化为墙钟节省。

\subsection{计算被移除的问题更新槽位}

对问题 $t$，令 $n_t$ 为它首次在已配置验证时点通过双边际验证器时的批量迭代编号，并令
\begin{equation}
A_k=\bigl|\{t:n_t\geq k\}\bigr|
\label{eq:survival-width}
\end{equation}
表示 packed update $k$ 之前的存活宽度。实测停止检查提供这两个量。

\begin{proposition}[问题更新槽位核算]
对于含有 $W$ 个候选的有限窗口，
\begin{equation}
S_{\mathrm{active}}
=\sum_{k=1}^{n_{\max}}A_k
=\sum_{t=1}^{W}n_t,
\qquad
S_{\mathrm{static}}=Wn_{\max}.
\label{eq:slot-identity}
\end{equation}
因此，相对于具有相同最终边界的固定宽度批次，动态压缩恰好移除
\begin{equation}
R=Wn_{\max}-\sum_{t=1}^{W}n_t
\label{eq:removable-slots}
\end{equation}
个候选更新槽位。
\end{proposition}
\begin{proof}
使用 $n_t=\sum_{k\geq1}\mathbf 1\{n_t\geq k\}$，并交换两个有限求和的顺序。固定宽度执行让全部 $W$ 条 lane 一直保持到 $n_{\max}$。
\end{proof}

逻辑 masking 会改变哪些状态发生更新，但不会改变 kernel 看到的宽度，因此仍保留 $Wn_{\max}$ 个槽位。静态矩形与生存曲线之间的面积，正是动态压缩在运行结束后实际移除的工作。检查间隔已经反映在观测到的 $n_t$ 中；在更密检查时点上首次通过的迭代编号会定义另一个量。

\begin{proposition}[任务无关的机会条件]
令 $(n_1,\ldots,n_W)$ 为任意有限迭代编号随机向量。那么 $\mathbb E[R]>0$ 当且仅当存在某次批量迭代 $k$ 满足
\begin{equation}
\Pr(0<A_k<W)>0.
\label{eq:joint-iteration-condition}
\end{equation}
\end{proposition}
\begin{proof}
$R$ 非负，并且仅当每条 lane 具有相同迭代编号时为零。不同的有限迭代编号会产生一个同时含有存活与已完成 lane 的 $k$ 值；反向结论直接由 $A_k$ 的定义得到。
\end{proof}

该条件取决于 packed window 内迭代编号的联合分布。即使不同窗口之间所需的迭代编号发生变化，完全同步的 lane 也不会产生可移除工作。

\subsection{实测后端上的精确墙钟核算}

令 $c_k(a,\xi_k)$ 为匹配更新 $k$ 在宽度 $a$ 时的成本，其中 $\xi_k$ 记录支撑形状、数据类型、kernel 配置以及对照中保持不变的其他后端状态。则
\begin{equation}
\begin{split}
T_{\mathrm{active}}={}&\sum_k c_k(A_k,\xi_k)
+H_{\mathrm{screen}}+H_{\mathrm{check}}\\
&+H_{\mathrm{compact}}+H_{\mathrm{other}},
\end{split}
\label{eq:time-accounting}
\end{equation}
$H_{\mathrm{other}}$ 包含其余计时内控制工作。令 $H_{\mathrm{fixed}}$ 为对应的固定宽度控制成本，并令 $\Delta H$ 表示 active 控制成本减去 $H_{\mathrm{fixed}}$。

\begin{proposition}[机会--转化分解]
对于具有相同最终验证边界和匹配更新状态 $\xi_k$ 的固定宽度/动态压缩对照，
\begin{equation}
T_{\mathrm{fixed}}-T_{\mathrm{active}}
=\sum_{k=1}^{n_{\max}}\left[c_k(W,\xi_k)-c_k(A_k,\xi_k)\right]-\Delta H.
\label{eq:time-difference}
\end{equation}
因此，动态压缩更快当且仅当
\begin{equation}
\sum_k\!\left[c_k(W,\xi_k)-c_k(A_k,\xi_k)\right]
>\Delta H.
\label{eq:active-gate}
\end{equation}
\end{proposition}
\begin{proof}
写出 $T_{\mathrm{fixed}}=\sum_k c_k(W,\xi_k)+H_{\mathrm{fixed}}$，减去式 \eqref{eq:time-accounting}，再合并控制项即可。
\end{proof}

这一分解呈现两道门槛。$n_t$ 之间的差异创造产生正可移除工作所需的存活宽度；实测函数 $c_k(a,\xi_k)$ 决定这些宽度缩减是否跨入更低的硬件成本层级。若所有 $k$ 都满足 $c_k(A_k,\xi_k)=c_k(W,\xi_k)$，动态压缩会减少问题更新槽位，却不会节省更新时间；若总成本下降超过 $\Delta H$，则产生净墙钟收益。

\subsection{波次形状的宽度成本}

对于本文使用的 packed Triton 更新，其 launch grid 含有
\begin{equation}
B(a,n)=a\left\lceil\frac{n}{b_m}\right\rceil
\label{eq:block-grid}
\end{equation}
个 thread block，其中 $b_m$ 是行 tile 大小。GPU 调度按受 occupancy 影响的波次处理该网格 \citep{yu2021habitat,nvidia2024cuda}。因此，一个有用的局部模型是
\begin{equation}
c(a;n,\xi)\approx g_{n,\xi}\!\left(\left\lceil\frac{B(a,n)}{Q_{\xi}}\right\rceil\right),
\label{eq:wave-cost}
\end{equation}
其中 $Q_{\xi}$ 是该 kernel 实测的驻留 block 容量，$g_{n,\xi}$ 把 wave 数映射为时间。寄存器、共享内存和 kernel 状态通过 $Q_{\xi}$ 进入模型；该数值针对所选设备与 kernel 配置进行校准。

在 $b_m=64$ 的实测 A100/Triton 配置上，第一波估计将饱和宽度分别放在 $n=1024$ 时的约 $6.75$、$n=4096$ 时的约 $1.69$，以及 $n=16384$ 时的一以下。直接原语测量在前两个规模上分别把首次大幅增长定位在宽度 $6\rightarrow7$ 与 $1\rightarrow2$；$n=16384$ 曲线从宽度一开始就对宽度变化作出响应。这些测量为该后端实例化 $c(a)$，并解释为何相同槽位缩减在不同问题规模上具有不同的墙钟价值。
